% Obsahový zdroj: import_hodiny/36_Kombinacni_cislo.tex a odpovídající PDF (obě strany).
% P1–P5: zdrojové Příklady 1–5; C1–C5: zdrojové Procvičování 1–5.
% Matematický obsah zachován; výklad, odvození a řešení pedagogicky rozvedeny.
% Upřesnění a opravy: viz KMB_29-41_Kontrola.md; původní import zůstává beze změny.
% Rozpis: 4_rocnik/tematicke_plany/Rozpis_vyuky_MAT_4_rocnik.tex, hodina 36.
% Plán: 4_rocnik/tematicke_plany/SPS.pdf, oddíly 2.1–2.3.
% Sazba: sablony/pripravy/TEMPLATE_Priprava.tex a společný priprava.sty.
% Vlastní doplnění: didaktický výklad, zdůvodnění, metodika; žádné nové číslované úlohy.
% Učebnice/09_Kombinatorika nemá obsahovou kapitolu; místní kombinatorické listy ani Realisticky.cz nejsou dostupné.
\documentclass[11pt,a4paper]{article}
% Spolecny styl pro ucitelske pripravy MAT (Work/MAT)
% Verze 1.1, 2026-10-23
\NeedsTeXFormat{LaTeX2e}
\ProvidesPackage{priprava}[2026/10/23 v1.1 Sablona ucitelskych priprav MAT]

% Zakladni typografie a cestina
\RequirePackage[T1]{fontenc}
\RequirePackage[utf8]{inputenc}
\RequirePackage[czech]{babel}
\RequirePackage{lmodern}
\RequirePackage{microtype}
\RequirePackage{amsmath,amssymb}
\RequirePackage{xcolor}
\RequirePackage{enumitem}
\RequirePackage[most]{tcolorbox}
\RequirePackage{tabularx,booktabs,array}
\RequirePackage{fancyhdr}
\RequirePackage{needspace}
\RequirePackage{geometry}
\RequirePackage[hidelinks]{hyperref}

\geometry{a4paper,top=17mm,bottom=18mm,left=23mm,right=23mm}
\definecolor{notegray}{gray}{0.39}
\setlength{\parindent}{0pt}
\setlength{\parskip}{2.5pt plus 1pt minus 0.5pt}
\setlist[itemize]{topsep=2pt,itemsep=1pt,parsep=0pt,leftmargin=1.45em}
\setlist[enumerate]{topsep=2pt,itemsep=2pt,parsep=0pt,leftmargin=1.65em}
\setcounter{secnumdepth}{0}
\renewcommand{\arraystretch}{1.12}

% Zahlavi: vlevo rocnik | globalni tema, vpravo cislo: nazev hodiny
% Zapati zustava zcela prazdne.
\pagestyle{fancy}
\fancyhf{}
\renewcommand{\headrulewidth}{0.35pt}
\renewcommand{\footrulewidth}{0pt}
\setlength{\headheight}{14pt}
\newcommand{\nastavPripravu}[4]{%
  \fancyhead[L]{\small #1\enspace\textbar\enspace #2}%
  \fancyhead[R]{\small #3: #4}%
}

% Minimalisticke cernobile boxy. Kazda definice zvlast.
\tcbset{commonlesson/.style={enhanced,breakable=false,colback=white,colframe=black,
  colbacktitle=white,coltitle=black,boxrule=0.45pt,
  arc=0.8mm,outer arc=0.8mm,left=3mm,right=3mm,top=1.6mm,bottom=1.6mm,
  toptitle=1mm,bottomtitle=0.7mm,fonttitle=\bfseries,
  attach boxed title to top left={xshift=2.5mm,yshift=-1.15mm},
  boxed title style={colback=white,colframe=white,boxrule=0pt,left=1mm,right=1mm}}}
\newtcolorbox{zadaniBox}[1][]{commonlesson,title={Zadání},#1}
\newtcolorbox{prikladBox}[1][]{commonlesson,title={Příklad},#1}
\newtcolorbox{definiceBox}[1][]{commonlesson,title={Definice},#1}

% Spolecny box pro zadani a jeho reseni, bez mezery mezi castmi.
% Pouziti:
%   \begin{prikladReseniBox}
%     Zadani ulohy.
%     \tcblower
%     \begin{align*} ... \end{align*}
%   \end{prikladReseniBox}
% Spodni cast ma jemny sedy podklad a delici linku; lze v ni sazet
% vice radku matematiky. Obe casti maji jen jeden spolecny ramecek.
\tcbset{solutionlesson/.style={
  commonlesson,
  enhanced jigsaw,
  breakable,
  colbacklower=black!3,
  segmentation style={solid,draw=black!25,line width=0.35pt},
  fontlower=\small\color{notegray},
  before lower={\noindent\textbf{Řešení:}\enspace},
  leftlower=3mm,rightlower=3mm
}}
\newtcolorbox{prikladReseniBox}[1][]{solutionlesson,title={Příklad},#1}
\newtcolorbox{zadaniReseniBox}[1][]{solutionlesson,title={Zadání},#1}

% Pomucky pro autora pripravy: sedive poznamky; \reseni pro starsi pripravy
\newcommand{\poznamka}[1]{\par\begingroup\small\color{notegray}\itshape
  \textnormal{\textbf{Poznámka pro učitele.}}\enspace #1\par\endgroup}
\newcommand{\reseni}[1]{\par\begingroup\small\color{notegray}
  \textbf{Očekávané řešení:}\enspace #1\par\endgroup}
\newcommand{\blok}[1]{\needspace{5\baselineskip}\subsection*{#1}}
\newcommand{\vyrok}[1]{\uv{#1}}

\nastavPripravu{4. ročník}{Kombinatorika}{36}{Kombinační číslo}
\begin{document}


Z předchozí hodiny již známe význam výběru skupiny. Nyní oddělíme tento význam od samotného početního zápisu: stejná kombinační čísla budeme počítat úpravou faktoriálů a používat v rovnicích.
\poznamka{Předpoklady: kombinace bez opakování, $0!$, krácení a kvadratické rovnice. Cíl: žák rychle vyčíslí kombinační číslo, vysvětlí krajní hodnoty a řeší rovnice pouze v přípustném celočíselném oboru. P5 je návrat ke slovní interpretaci výsledku.}

\blok{Výklad a odvození}
\begin{definiceBox}
\textbf{Kombinační číslo} $\binom nk$, čtené „$n$ nad $k$“, pro celá $n\geq0$ a $0\leq k\leq n$ definujeme jako
\[
 \binom nk=\frac{n!}{k!(n-k)!}.
\]
Udává počet $k$-prvkových podmnožin $n$-prvkové množiny. Spodní index $k$ je rovněž celé číslo.
\end{definiceBox}

Z výběrového významu i dosazením $0!=1$ plyne
\[
 \binom n0=\binom nn=1\quad(n\geq0),\qquad
 \binom n1=n\quad(n\geq1).
\]
Pro celé $n\geq2$ navíc
\[
 \binom n2=\frac{n(n-1)(n-2)!}{2!(n-2)!}=\frac{n(n-1)}2.
\]
Prázdný výběr je jedna možnost, ne nula možností. Výsledek je nezáporné celé číslo, protože počítá skupiny.

Než budeš násobit, vyhledej společný faktoriál. V $\binom{12}{10}$ lze krátit $10!$, takže zůstanou jen dva činitelé v čitateli a $2!$ ve jmenovateli. Obecný důkaz symetrie ponecháme do hodiny 37.

U rovnic nejprve určujeme obor: z $\binom n2$ vyplývá celé $n\geq2$. Teprve potom nahradíme kombinační číslo výrazem $n(n-1)/2$. Algebraická rovnice může mít i záporný kořen; takový kořen není řešením původní úlohy.


\poznamka{Průběh: 0–8 min význam a krajní hodnoty; 8–19 min P1–P2; 19–31 min P3–P4; 31–41 min C1–C5; 41–45 min kontrola a krátké slovní P5. Při pomalejším tempu ponech P5 jako domácí opakování hodiny 35.}

\blok{Společně řešené příklady}

% Zdroj P1: import_hodiny/36_Kombinacni_cislo.tex, Příklad 1.
\begin{prikladReseniBox}[title={Příklad P1}]
Vypočti \(\binom{10}{3}\).
\tcblower
\(\frac{10!}{3!7!}=\frac{10\cdot9\cdot8}{3\cdot2\cdot1}=120\).
\end{prikladReseniBox}


% Zdroj P2: import_hodiny/36_Kombinacni_cislo.tex, Příklad 2.
\begin{prikladReseniBox}[title={Příklad P2}]
Vypočti \(\binom{12}{10}\) bez velkých faktoriálů.
\tcblower
\(\binom{12}{10}=\frac{12!}{10!2!}=\frac{12\cdot11}{2}=66\).
\end{prikladReseniBox}


% Zdroj P3: import_hodiny/36_Kombinacni_cislo.tex, Příklad 3.
\begin{prikladReseniBox}[title={Příklad P3},breakable=false]
Urči celé $n$, pro které platí $\binom n2=45$.
\tcblower
Podmínka: celé $n\geq2$. Potom
\begin{align*}
 \binom n2=45&\iff n(n-1)=90\\
 &\iff n^2-n-90=0\\
 &\iff(n-10)(n+9)=0.
\end{align*}
Kořeny jsou $10$ a $-9$, druhý nevyhovuje podmínce. Zkouška: $\binom{10}{2}=10\cdot9/2=45$. Jediné řešení je $n=10$.
\end{prikladReseniBox}


% Zdroj P4: import_hodiny/36_Kombinacni_cislo.tex, Příklad 4.
\begin{prikladReseniBox}[title={Příklad P4}]
Vyřeš $\binom n2=3n$ v celých číslech $n\geq2$.
\tcblower
Pro celé $n\geq2$ je $n>0$, takže smíme dělit $n$:
\[
 \frac{n(n-1)}2=3n\iff n(n-1)=6n\iff n-1=6\iff n=7.
\]
Zkouška: $\binom72=21=3\cdot7$. Nula, která by vyšla po vynechání podmínky, nemá v tomto zadání význam.
\end{prikladReseniBox}


% Zdroj P5: import_hodiny/36_Kombinacni_cislo.tex, Příklad 5.
\begin{prikladReseniBox}[title={Příklad P5}]
Z osmi lidí se vybírá čtyřčlenná komise. Dva konkrétní kandidáti musejí oba být vybráni.
\tcblower
Dva konkrétní lidé jsou pevně vybráni, zbývají dvě volná místa. Z ostatních šesti kandidátů vybereme neuspořádanou dvojici, tedy $\binom62=15$ komisí. Nemůžeme vybírat další dva z osmi, neboť předepsané lidi nelze vybrat podruhé.
\end{prikladReseniBox}


\newpage
\blok{Samostatná práce a zásoba k procvičení}

C1–C2 počítej krácením. U C3–C4 napiš definiční obor a proveď zkoušku. U C5 vysvětli, co je jediným výsledkem každého z obou výběrů.

% Zdroj C1: import_hodiny/36_Kombinacni_cislo.tex, Procvičování, položka 1.
\begin{zadaniReseniBox}[title={Úloha C1}]
Spočti \(\binom83\).
\tcblower
Po zkrácení $5!$ dostáváme $\binom83=\dfrac{8\cdot7\cdot6}{3\cdot2\cdot1}=56$.
\end{zadaniReseniBox}

% Zdroj C2: import_hodiny/36_Kombinacni_cislo.tex, Procvičování, položka 2.
\begin{zadaniReseniBox}[title={Úloha C2}]
Spočti \(\binom{11}{9}\).
\tcblower
$\binom{11}{9}=\dfrac{11!}{9!\,2!}=\dfrac{11\cdot10}{2}=55$.
\end{zadaniReseniBox}

% Zdroj C3: import_hodiny/36_Kombinacni_cislo.tex, Procvičování, položka 3.
\begin{zadaniReseniBox}[title={Úloha C3}]
Najdi celé $n$, pro něž $\binom n1=13$.
\tcblower
Podmínka: celé $n\geq1$. Protože $\binom n1=n$, musí být $n=13$; zkouška $\binom{13}{1}=13$.
\end{zadaniReseniBox}

% Zdroj C4: import_hodiny/36_Kombinacni_cislo.tex, Procvičování, položka 4.
\begin{zadaniReseniBox}[title={Úloha C4}]
Najdi celé $n\geq2$, pro něž $\binom n2=28$.
\tcblower
Podmínka: celé $n\geq2$. Rovnice vede na $n(n-1)=56$, tedy $(n-8)(n+7)=0$. Z kořenů $8,-7$ vyhovuje jen $n=8$. Zkouška: $\binom82=28$.
\end{zadaniReseniBox}

% Zdroj C5: import_hodiny/36_Kombinacni_cislo.tex, Procvičování, položka 5.
\begin{zadaniReseniBox}[title={Úloha C5}]
Kolik možností má výběr všech sedmi prvků ze sedmiprvkové množiny? Kolik možností má výběr žádného prvku?
\tcblower
Při výběru všech prvků je jediným výsledkem celá množina; při výběru žádného jediným výsledkem prázdná množina. Proto $\binom77=\binom70=1$. Také oba faktoriálové zlomky mají hodnotu $7!/(7!\,0!)=1$.
\end{zadaniReseniBox}


\poznamka{Typické chyby: $\binom n0=0$; přijetí záporného kořene; dělení $n$ bez ověření $n\ne0$. Ve výpočtu $\binom{12}{10}$ není třeba počítat deset činitelů ani používat ještě neprobranou větu. Příště zdůvodníme symetrii, Pascalův vztah a součet všech podmnožin.}
\end{document}
